%==============================================================================
%  Borhaneddine Hamadou -- Curriculum Vitae
%
%  Compile with pdfLaTeX on Overleaf (Menu > Compiler > pdfLaTeX).
%  No shell-escape, no external images, no bibliography pass: one click builds it.
%
%  ATS notes (why the file looks the way it does):
%    * Single column. No text boxes, no multicol, no tikz, no images, no tables.
%    * Every "left text ....... right date" row is plain text plus \hfill, so a
%      parser reads the line exactly as printed, in order.
%    * Standard section names (Education, Professional Experience, Skills, ...).
%    * Icons are optional font glyphs placed NEXT TO the readable text, so the
%      email/phone/URL still parse correctly if a parser drops the glyph.
%
%  Two switches:
%    1. Icons off      -> change \showiconstrue to \showiconsfalse below.
%    2. Arial/Calibri  -> comment out \usepackage{XCharter} and uncomment the
%                         two \helvet lines just under it.
%
%  Tailoring per application (recommended -- a generic CV is easy to spot):
%    edit the headline, the "Seeking..." line, and the Summary so they mirror
%    the job ad's own wording (validation & verification, SOTIF, simulation,
%    embedded, ADAS, ...). Leave everything else fixed.
%==============================================================================

\documentclass[a4paper,10pt]{article}

%------------------------------ ENCODING & FONT -------------------------------
\usepackage[utf8]{inputenc}
\usepackage[T1]{fontenc}
\usepackage{textcomp}
\usepackage{microtype}

\usepackage{XCharter}                  % serif face matching the reference layout
% --- Arial / Calibri-like alternative: comment the line above, uncomment these:
% \usepackage{helvet}
% \renewcommand{\familydefault}{\sfdefault}

%--------------------------------- GEOMETRY -----------------------------------
\usepackage[a4paper,top=1.4cm,bottom=1.4cm,left=1.5cm,right=1.5cm,footskip=0.6cm]{geometry}

%---------------------------------- COLOUR ------------------------------------
\usepackage{xcolor}
\definecolor{accent}{HTML}{2E5C8A}

%----------------------------------- LINKS ------------------------------------
\usepackage{hyperref}
\hypersetup{
  colorlinks=true,
  urlcolor=accent,
  linkcolor=accent,
  pdftitle={Borhaneddine Hamadou - Curriculum Vitae},
  pdfauthor={Borhaneddine Hamadou},
  pdfsubject={Research Engineer / R\&D - Safety Validation and Testing of Autonomous Systems},
  pdfkeywords={autonomous systems, scenario-based testing, safety validation,
               verification and validation, SOTIF, ISO 21448, ISO 34503,
               Bayesian networks, simulation, CARLA, ROS 2, Python, R\&D}
}

%----------------------------------- ICONS ------------------------------------
\newif\ifshowicons
\showiconstrue                 % <-- set to \showiconsfalse to strip every glyph

\ifshowicons
  \usepackage{fontawesome5}
  \newcommand{\cvicon}[1]{{\small#1}~}
\else
  \newcommand{\cvicon}[1]{}
\fi

%--------------------------------- SECTIONS -----------------------------------
\usepackage{titlesec}
\titleformat{\section}
  {\color{accent}\large\bfseries}
  {}{0pt}{}[\vspace{-6pt}\color{accent}\rule{\textwidth}{0.9pt}]
\titlespacing*{\section}{0pt}{9pt}{5pt}

%----------------------------------- LISTS ------------------------------------
\usepackage{enumitem}
\setlist[itemize]{
  topsep=1pt, partopsep=0pt, parsep=0pt, itemsep=1.5pt,
  leftmargin=13pt, label={\small\textbullet}
}

\pagestyle{empty}
\setlength{\parindent}{0pt}
\raggedbottom

%------------------------------- ENTRY MACROS ---------------------------------
% \parfillskip=0pt inside the group makes \hfill the only stretchable glue on
% the line, so the right-hand text lands flush right instead of mid-line.

% \cventry{Organisation}{Location}{Role}{Dates}
\newcommand{\cventry}[4]{%
  \vspace{3pt}%
  \begingroup\parfillskip=0pt\relax
    \textbf{#1}\hfill #2\par
    \textit{#3}\hfill \textit{#4}\par
  \endgroup
  \vspace{1pt}%
}

% \cvproject{Title}{Stack / context}
\newcommand{\cvproject}[2]{%
  \vspace{3pt}%
  \begingroup\parfillskip=0pt\relax
    \textbf{#1}\hfill \textit{#2}\par
  \endgroup
  \vspace{1pt}%
}

% \cvline{Left text}{Right text} -- compact one-line entry
\newcommand{\cvline}[2]{%
  \begingroup\parfillskip=0pt\relax
    #1\hfill #2\par
  \endgroup
  \vspace{1pt}%
}

% \cvskill{Category}{Items}
\newcommand{\cvskill}[2]{%
  \textbf{#1}\hspace{0.6em}#2\par\vspace{2pt}%
}

\newcommand{\cvsep}{\hspace{4pt}\textbar\hspace{4pt}}
\newcommand{\cvdot}{\,\textperiodcentered\,}

%==============================================================================
\begin{document}

%---------------------------------- HEADER ------------------------------------
\begin{center}
  {\LARGE\bfseries\color{accent} Borhaneddine Hamadou}\\[4pt]
  {\large Research Engineer (R\&D) \textbar{} Safety Validation \& Testing of Autonomous Systems}\\[3pt]
  {\small Seeking an industrial R\&D / research engineer position in autonomous and safety-critical systems}\\[5pt]
  {\small
    \cvicon{\faPhone}+33 7 84 34 66 86 \cvsep
    \cvicon{\faEnvelope}\href{mailto:hamadou.borhaneddine@gmail.com}{hamadou.borhaneddine@gmail.com} \cvsep
    \cvicon{\faLinkedin}\href{https://www.linkedin.com/in/borhaneddine-hamadou-48272431b/}{linkedin.com/in/borhaneddine-hamadou} \cvsep
    \cvicon{\faGithub}\href{https://github.com/BorhaneddineHamadou}{github.com/BorhaneddineHamadou}\\[2pt]
    \cvicon{\faGlobe}\href{https://borhaneddinehamadou.github.io/}{borhaneddinehamadou.github.io} \cvsep
    \cvicon{\faMapMarkerAlt}\^Ile-de-France, France
  }
\end{center}
\vspace{-2pt}

%---------------------------------- SUMMARY -----------------------------------
\section*{Summary}

PhD researcher in \textbf{safety validation and testing of autonomous driving systems} at CEA List (Universit\'e Paris-Saclay), defending July 2027. Three years building testing methods --- scenario generation, test adequacy metrics, regression test selection --- grounded in \textbf{ISO 21448 (SOTIF)} and \textbf{ISO 34503} and validated on production driving stacks (Autoware, openpilot) in \textbf{CARLA} and \textbf{MetaDrive}. Published at \textbf{ASE} and \textbf{ISSRE}, and shipped as open-source tooling. Previously machine learning for autonomous-vehicle edge computing at \textbf{LAAS-CNRS}.

%--------------------------------- EDUCATION ----------------------------------
\section*{Education}

\cventry{Universit\'e Paris-Saclay --- CEA List}{Palaiseau, France}
      {PhD in Computer Science --- Safety Testing of Autonomous Vehicle Systems}{Jul 2024 --- Jul 2027 (expected)}
\begin{itemize}
  \item Thesis on scenario-based testing and safety validation of autonomous driving systems, combining Bayesian networks, combinatorial coverage, and statistical modelling.
  \item Methods evaluated on open-source production stacks (Autoware, openpilot) in CARLA and MetaDrive simulation, against ISO 21448 (SOTIF) and ISO 34503 requirements.
\end{itemize}

\cventry{\'Ecole Nationale Sup\'erieure d'Informatique (ESI)}{Algiers, Algeria}
      {Master of Engineering (Ing\'enieur d'\'Etat) in Computer Systems Engineering}{Sep 2018 --- Sep 2023}
\begin{itemize}
  \item Leading computer science school in Algeria. Specialisation in machine learning, combinatorial optimisation, data analysis, high-performance computing, distributed systems, and network and systems security.
  \item Engineering thesis (PFE) carried out at LAAS-CNRS, Toulouse --- see Professional Experience.
\end{itemize}

%------------------------------- EXPERIENCE -----------------------------------
\section*{Professional Experience}

\cventry{CEA List --- Universit\'e Paris-Saclay}{Palaiseau, France}
      {PhD Researcher (R\&D) --- Autonomous Systems Testing}{Jul 2024 --- Present}
\begin{itemize}
  \item Designed \textbf{PWCC}, a risk-informed test adequacy metric weighting combinatorial coverage by operational probability learned from fleet data: the \textbf{strongest predictor of fault detection} across 4 autonomous driving datasets and 7 baselines (Spearman $r_s = 0.981$), removing the false ``adequate'' verdicts behind which standard coverage hides up to \textbf{95.8\%} of failure-triggering conditions. \emph{Published at ASE 2026.}
  \item Built \textbf{BayScen}, an ISO 34503-grounded scenario generator pairing capability abstraction with Bayesian networks learned from real-world data: critical scenarios \textbf{$\times$7 more likely to be physically plausible} than adversarial search, at the highest ODD coverage among the methods compared. \emph{Published at ISSRE 2026.}
  \item Developed \textbf{SA-TRP}, a stochasticity-aware regression test selection method for non-deterministic autonomous driving stacks, improving regression-detection F1 by \textbf{34--62\%} on openpilot and Autoware commit histories. \emph{Under review at ICSE 2027.}
  \item Released PWCC as a peer-reviewed replication package --- CI-ready command-line tool, Docker image, 31 automated tests, fully pinned dependencies --- archived on Zenodo under CC BY 4.0 and reproducible end to end in under 30 minutes.
  \item Implemented the full experimental pipeline in Python (pgmpy, scikit-learn, NumPy, pandas) over large-scale simulation campaigns; presented results to academic and industrial audiences in France and Brazil.
\end{itemize}

\cventry{LAAS-CNRS}{Toulouse, France}
      {AI Research Intern --- Autonomous Vehicles \& Edge Computing}{Oct 2022 --- Sep 2023}
\begin{itemize}
  \item Engineered time-series forecasting models predicting Quality-of-Service degradation for autonomous vehicles in dynamic edge-computing environments, enabling intervention \emph{before} latency incidents occur.
  \item Delivered the models as a real-time prediction REST API, integrated into a larger autonomous-vehicle ecosystem for direct deployment.
\end{itemize}

\cventry{Universit\'e Paris-Saclay}{France}
      {Teaching Assistant --- Computer Science}{Jan 2025 --- Present}
\begin{itemize}
  \item Laboratory and tutorial sessions in Object-Oriented Programming (Java) and Database Management Systems (Oracle: relational design, SQL, PL/SQL, administration), including supervision of student projects.
\end{itemize}

\vspace{3pt}
\cvline{\small\textbf{Software Developer (Freelance)} --- Somu Raqami, Algeria \cvdot Full-stack mobile \& web}{\small Apr 2019 --- Sep 2022}
\cvline{\small\textbf{Software Developer Intern} --- Maystro Delivery, Algiers, Algeria \cvdot Stock, sales \& POS web app}{\small Jul --- Sep 2021}

%--------------------------------- PROJECTS -----------------------------------
\section*{Selected Projects}

\cvproject{PWCC --- Test Adequacy Tool for Autonomous Systems}{Python, pgmpy, Docker, pytest \cvdot CEA List, 2025--2026}
\begin{itemize}
  \item Open-source command-line tool that learns an operating-domain model from a user's own scenario data, reports which operating conditions still need testing and how much exposure each requires, and exits non-zero when testing is inadequate --- usable directly as a CI stopping gate. \href{https://doi.org/10.5281/zenodo.21739211}{doi.org/10.5281/zenodo.21739211}
\end{itemize}

\cvproject{QoS Prediction API for Autonomous Vehicles}{Python, scikit-learn, REST \cvdot LAAS-CNRS, 2022--2023}
\begin{itemize}
  \item End-to-end machine-learning pipeline and production-style prediction service for Quality-of-Service degradation in autonomous-vehicle edge environments. \href{https://github.com/BorhaneddineHamadou/QoS_Pred_API_V2}{github.com/BorhaneddineHamadou/QoS\_Pred\_API\_V2}
\end{itemize}

\cvproject{High-Performance Solver \& Hybrid Metaheuristics}{C, CUDA, OpenMP, MPI \cvdot ESI Algiers}
\begin{itemize}
  \item Sequential and parallel CUDA/OpenMP Gauss-Jordan linear solvers benchmarked under high-computation loads; exact, heuristic, metaheuristic, and hybrid solvers for the unbounded knapsack problem.
\end{itemize}

% \cvproject{SpeedBoat Industry --- Strategic Decision Dashboard}{JavaScript, D3.js, Chart.js, Three.js \cvdot CentraleSup\'elec, 2026}
% \begin{itemize}
%   \item Interactive executive dashboard (HR, sales and incident analytics, risk matrix, life-cycle analysis, 3D model) built in three days with a team of four PhD students during the ``Float and Serious\textregistered'' doctoral hackathon.
% \end{itemize}

%------------------------------- PUBLICATIONS ---------------------------------
\section*{Publications}

\begin{itemize}
  \item \textbf{B. Hamadou}, C. Mraidha, N. Matragkas. ``PWCC: Probability-Weighted Combinatorial Coverage for Risk-Informed Scenario Testing of Autonomous Systems.'' \emph{ASE 2026 --- 41st IEEE/ACM International Conference on Automated Software Engineering (Research Track)}, Munich, Germany.
  \item \textbf{B. Hamadou}, C. Mraidha, N. Matragkas. ``BayScen: Standards-Grounded Scenario Generation for Autonomous Vehicle Testing via Bayesian Networks and Capability Abstraction.'' \emph{ISSRE 2026 --- IEEE International Symposium on Software Reliability Engineering (Research Track)}.
  \item \textbf{B. Hamadou}, C. Mraidha, N. Matragkas. ``Stochasticity-Aware Test Reduction and Prioritization for Autonomous Driving Systems.'' \emph{Under review at ICSE 2027}.
  \item \textbf{B. Hamadou}. ``Bayesian Network-Based Scenario Generation for Autonomous Systems Testing.'' \emph{ISSRE 2025 --- IEEE International Symposium on Software Reliability Engineering (Doctoral Symposium)}, S\~ao Paulo, Brazil.
\end{itemize}

%-------------------------- CERTIFICATIONS & TRAINING -------------------------
\section*{Certifications \& Training}

\begin{itemize}
  \item \textbf{ISTQB Certified Tester Foundation Level} --- training course, G2C/GTEC, Tunis, Tunisia (Oct 2024). Test process and test levels, static testing, black-box and white-box test design techniques, test management.
  \item \textbf{Parcours DRIVE\textregistered{} --- R\&D en Entreprise}, Universit\'e Paris-Saclay (2025--2026). Doctoral programme on industrial R\&D careers, concluded by a three-day corporate-strategy hackathon at CentraleSup\'elec.
  \item \textbf{ESSAI 2025 --- European Summer School on Artificial Intelligence}, Bratislava, Slovakia (Jul 2025).
  \item \textbf{Formation Ing\'enieur Entreprendre} --- entrepreneurship diploma, Maison de l'Entrepreneuriat, ESI Algiers (2022--2023).
\end{itemize}

%---------------------------------- SKILLS ------------------------------------
\section*{Technical Skills}

\cvskill{Testing \& Validation}{Scenario-based testing \cvdot Combinatorial testing \cvdot Test adequacy \& coverage metrics \cvdot Regression test selection \cvdot Verification \& validation (V\&V) \cvdot ISO 21448 (SOTIF) \cvdot ISO 34503 (ODD) \cvdot ISTQB foundations}
\cvskill{Simulation \& Robotics}{CARLA \cvdot MetaDrive \cvdot Gazebo \cvdot ROS 2 \cvdot Autoware \cvdot openpilot}
\cvskill{ML \& Statistics}{Bayesian networks (pgmpy) \cvdot Probabilistic modelling \cvdot Monte Carlo methods \cvdot Time-series forecasting \cvdot Supervised \& unsupervised learning \cvdot scikit-learn \cvdot TensorFlow \cvdot NumPy \cvdot pandas}
\cvskill{Programming}{Python \cvdot C \cvdot C++ \cvdot Java \cvdot SQL / PL/SQL \cvdot Bash \cvdot JavaScript}
\cvskill{HPC \& Parallel}{CUDA \cvdot OpenMP \cvdot MPI \cvdot Pthreads}
\cvskill{Tooling}{Linux \cvdot Git \cvdot Docker \cvdot pytest \cvdot Continuous integration \cvdot \LaTeX}

%--------------------------- AWARDS & ACADEMIC SERVICE ------------------------
\section*{Awards \& Academic Service}

\begin{itemize}
  \item \textbf{Best Product Award (``Product of the Year'')} --- INJAZ Algeria national competition for young entrepreneurs, 2022.
  \item \textbf{5th place} --- Algerian Collegiate Programming Contest (ICPC Foundation), 2022.
  \item \textbf{Organising committee}, JITA 2025 --- Junior Conference on Informatics: Theory and Applications, Universit\'e Paris-Saclay: communication, visual design, website, and review process.
  % \item \textbf{Presenter}, ISSRE 2025 Doctoral Symposium, S\~ao Paulo, Brazil; annual thesis-progress talks at CEA PhD Day (2025, 2026).
\end{itemize}

%--------------------------------- LANGUAGES ----------------------------------
\section*{Languages}

\textbf{Arabic} (native) \cvsep \textbf{French} (full professional proficiency --- language of my teaching and of regular research presentations) \cvsep \textbf{English} (full professional proficiency --- research published and presented internationally) \cvsep \textbf{German} (A2)

\vspace{7pt}
{\small Available from July 2027 \cvdot Open to relocation across France and Europe \cvdot References available upon request.}

\end{document}
